	\chapter{Differential Calculus}
	\section{Basic Concepts}
	\begin{definitionenv}
		Let $f$ be a function. If
		\[
		\lim _{h\to 0} \frac{f(x+h)-f(x)}{h}=L \in \mathbb{R},
		\]
		then we call it the \textbf{derivative} of $f$ and denote it by	$f'(x)$. Furthermore, the $n$-th derivative of $f$ is defined by
		\[
		f^{(n)}(x)= \left[f^{(n-1)}(x)\right]' \quad \text{for}\quad n=1,2,3,\cdots
		\]
		and $f^{(0)}(x)=f(x)$.
	\end{definitionenv}

	\begin{exampleenv}
		We're going to prove some basic derivatives using the definition.
		\begin{enumerate}
			\item Suppose $n$ is a positive integer and $x\ge 0$. Consider $f(x)=x^n$, then
			\[
			\begin{aligned}
				f'(x) &= \lim_{h\to 0} \frac{(x+h)^n - x^n}{h}\\ &= \lim_{h\to 0} \sum_{k=0}^{n-1} (x+h)^k x^{n-1-k}\\ &= \sum_{k=0}^{n-1} x^k x^{n-1-k}\\ &= n x^{n-1}.
			\end{aligned}
			\]
			\item Consider $f(x)=\sin x$, then
			\[
			\begin{aligned}
				f'(x) &= \lim_{h\to 0} \frac{\sin(x+h) - \sin x}{h}\\[6pt] 
				&= \lim_{h\to 0} \dfrac{2\sin\left(\dfrac{h}{2}\right)\cos\left(x+\dfrac{h}{2}\right)}{h}\\[6pt] 
				&= \left[ \lim_{h \to 0} \dfrac{\sin\left(\dfrac{h}{2}\right)}{\dfrac{h}{2}}\right] \left[\lim_{h \to 0} \cos\left(x+\dfrac{h}{2}\right)\right]  \\ 
				&= \cos x.
			\end{aligned}
			\]
			\item Consider $f(x)=\cos x$, then
			\[
			\begin{aligned}
				f'(x) &= \lim_{h\to 0} \frac{\cos(x+h) - \cos x}{h}\\[6pt] 
				&= \lim_{h\to 0} \dfrac{-2\sin\left(\dfrac{h}{2}\right)\sin\left(x+\dfrac{h}{2}\right)}{h}\\[6pt] 
				&= -\left[ \lim_{h \to 0} \dfrac{\sin\left(\dfrac{h}{2}\right)}{\dfrac{h}{2}}\right] \left[\lim_{h \to 0} \sin\left(x+\dfrac{h}{2}\right)\right]  \\ 
				&= -\sin x.
			\end{aligned}
			\]
			\item Suppose $n$ is a positive integer and $x>0$. Consider $f(x)=x^{\frac{1}{n}}$ and for convenience let 
			\[
			u=\left( x+h \right)^{\frac{1}{n}} \quad\text{and}\quad v=x^{\frac{1}{n}}.
			\]
			Hence we have
			\[
			\begin{aligned}
				f'(x) &= \lim_{h\to 0} \frac{(x+h)^{\frac{1}{n}} - x^{\frac{1}{n}}}{h}\\[6pt] 
				&= \lim_{h\to 0} \dfrac{u - v}{u^n - v^n} \\[6pt] 
				&= \lim_{u\to v} \dfrac{1}{u^{n-1} + u^{n-2}v + \cdots + uv^{n-2} + v^{n-1}} \\[6pt]
				&= \dfrac{1}{n x^{\frac{n-1}{n}}}\\[6pt]
				&= \dfrac{1}{n} x^{\frac{1}{n}-1}.
			\end{aligned}
			\]
			For the case $x=0$, we can see $f(x)$ is defined at $x=0$ but $f'(x)$ is not defined.
		\end{enumerate}
	\end{exampleenv}
	\begin{theoremenv}
		If function $f$ has a derivative at $x$, then $f$ is continuous at $x$.
	\end{theoremenv}
	\begin{proofenv}
		Let
		\[
		f'(x)=\lim_{h\to 0} \frac{f(x+h)-f(x)}{h}.
		\]
		So
		\[
		\lim_{h\to 0} f(x+h) = \lim_{h\to 0} \left[ \frac{f(x+h)-f(x)}{h} \cdot h + f(x) \right] = f'(x)\cdot 0 + f(x) = f(x).
		\]
	\end{proofenv}
    \begin{remarkenv}
        The converse of the above theorem is not true, i.e., if $f$ is continuous at $x$, then $f$ does not necessarily have a derivative at $x$. The absolute value function $f(x)=|x|$ is a typical counterexample since $f$ is continuous at $0$ but $f'(0)$ does not exist.

        Moreover, even if $f$ is differentiable at $x$, $f'$ may not be continuous at $x$. For example, we define
        $$
        f(x)=
        \begin{cases}
        x^2\sin\dfrac{1}{x^2}, & x\ne 0,\\[6pt]
        0, & x=0.
        \end{cases},\qquad
        f'(x)= \begin{cases}
            2x\sin\dfrac{1}{x^2}-\dfrac{2}{x}\cos\dfrac{1}{x^2}, & x\ne 0,\\[6pt]
            0, & x=0.
        \end{cases}.
        $$
        Then $f$ is differentiable on $\R$ but $f'$ is not continuous at $0$ since $\displaystyle\lim_{x\to 0} f'(x)$ does not exist while $f'(0)=0$.
    \end{remarkenv}
	\begin{theoremenv}
		Suppose $f$ and $g$ are two functions defined on a common interval. If both $f$ adn $g$ have derivatives at $x=a$, then so do the following functions:
		\[
		f+g,f-g,fg, \frac{f}{g}\ \text{(provided that $g(a)\neq 0$)}.
		\]
		Additionally, their derivatives satisfy the following:
		\begin{enumerate}
			\item[(1)] $(f+g)'=f'+g'$,
			\item[(2)] $(f-g)'=f'-g'$,
			\item[(3)] $(fg)'=f'g+fg'$,
			\item[(4)] $\left(\dfrac{f}{g}\right)'=\dfrac{f'g-fg'}{g^2}$.
		\end{enumerate}
	\end{theoremenv}
	\begin{proofenv}
		We only prove (3) and (4) here.
		\begin{enumerate}
			\item[(3)] We have
			\[
			\begin{aligned}
				(fg)'(a) &= \lim_{h\to 0} \frac{f(a+h)g(a+h) - f(a)g(a)}{h}\\[6pt]
				&= \lim_{h\to 0} \left[ g(a+h)\cdot \frac{f(a+h)-f(a)}{h} + f(a)\cdot \frac{g(a+h)-g(a)}{h} \right]\\[6pt]
				&= g(a)\cdot f'(a) + f(a)\cdot g'(a).
			\end{aligned}
			\]
			\item[(4)] We have
			\[
			\begin{aligned}
				\left(\frac{f}{g}\right)'(a) &= \lim_{h\to 0} \frac{\dfrac{f(a+h)}{g(a+h)} - \dfrac{f(a)}{g(a)}}{h}\\[6pt]
				&= \lim_{h\to 0} \frac{f(a+h)g(a) - f(a)g(a+h)}{h g(a+h) g(a)}\\[6pt]
				&= \lim_{h\to 0} \left[ \frac{g(a)\cdot [f(a+h)-f(a)]}{h g(a+h) g(a)} - \frac{f(a)\cdot [g(a+h)-g(a)]}{h g(a+h) g(a)} \right]\\[6pt]
				&= \frac{g(a)\cdot f'(a)}{g^2(a)} - \frac{f(a)\cdot g'(a)}{g^2(a)}\\[6pt]
				&= \frac{f'g - fg'}{g^2}\bigg|_{x=a}.
			\end{aligned}
			\]
		\end{enumerate}
	\end{proofenv}
	\begin{remarkenv}
		Using these properties and induction, we can expand
		\[
		\left( x^n\right)' = n x^{n-1}
		\]
		for all rational numbers $n$.
	\end{remarkenv}
	\begin{notationenv}
		In fact, the following notations are also commonly used for derivatives:
		\begin{enumerate}
			\item Leibniz's notation:
			\[
			\frac{\dd y}{\dd x},\ \frac{\dd^2 y}{\dd x^2},\ \frac{\dd^3 y}{\dd x^3},\ \cdots,\ \frac{\dd^n y}{\dd x^n},
			\]
			where $\dd x$ and $\dd y$ are called \textbf{differentials}.
			\item Lagrange's notation:
			\[
			f',\ f'',\ f''',\cdots,\ f^{(n)};
			\]
			\item Newton's notation (mainly used in physics):
			\[
			\dot{y},\ \ddot{y},\ \dddot{y},\ \cdots;
			\]
			\item Arbogast's notation:
			\[
			Df,\ D^2f,\ D^3f,\ \cdots,\ D^n f,
			\]
			where $D$ is called the \textbf{differential operator}. It's easy to verify that $D$ is linear.
		\end{enumerate}
	\end{notationenv}

	\section{Chain Rule for Differentiation}
	\begin{theoremenv}
		Suppose $u$ and $v$ are functions and the function $f$ is given by
		$f=u\circ v$. Suppose $v'(x)$ and $u'(y)$ where $y=v(x)$ both exist. Then $f'(x)$ exists and is given by
		\[
		f'(x) = u'(v(x)) \cdot v'(x).
		\]
	\end{theoremenv}
	\begin{proofenv}
		For convenience, let
		\[
		y=v(x),\quad \Delta y = v(x+h) - v(x).
		\]
		% Then if we assume $\Delta y\neq 0$,  we have
		% \[
		% \begin{aligned}
		% 	f'(x) &= \lim_{h\to 0} \frac{u(v(x+h)) - u(v(x))}{h}\\[6pt]
		% 	&= \lim_{h\to 0} \frac{u(y + \Delta y) - u(y)}{h}\\[6pt]
		% 	&= \lim_{h\to 0} \left[ \frac{u(y + \Delta y) - u(y)}{\Delta y} \cdot \frac{\Delta y}{h} \right]\\[6pt]
		% 	&= \left[ \lim_{\Delta y \to 0} \frac{u(y + \Delta y) - u(y)}{\Delta y} \right] \cdot \left[ \lim_{h\to 0} \frac{\Delta y}{h} \right]\\[6pt]
		% 	&= u'(v(x)) \cdot v'(x).
		% \end{aligned}
		% \]
		Let
		\[
		g(t)= 
		\begin{cases}
		0, & \text{if}\ t=0,\\[6pt]
		\dfrac{u(y+t) - u(y)}{t}-u'(y), & \text{if}\ t\neq 0,
		\end{cases}
		\]
		since we know for sure
		\[
		\lim_{t\to 0} g(t) = \lim_{t\to 0} \left[ \frac{u(y+t) - u(y)}{t} - u'(y) \right]=0.
		\]
		Then we have
		\[
		u(y+t)-u(y)=
		\begin{cases}
		0\left[u'(y) + g(t)\right], & \text{if}\ t=0\\[6pt]
		t\left[u'(y) + g(t)\right], & \text{if}\ t\neq 0
		\end{cases}
		= t\left[u'(y) + g(t)\right]
		\]
		no matter $t=0$ or not. Thus plugging in $t=\Delta y$, we know
		\[
		f'(x)=\lim_{h\to 0} \frac{u(y + \Delta y) - u(y)}{h} = \lim_{h\to 0} \frac{\Delta y}{h} \left[u'(y) + g(\Delta y)\right].
		\]
		Note that
		\[
		\lim_{h\to 0} \frac{\Delta y}{h}=v'(x)
		\]
		as discussed before. Also we know $g$ is continuous at $0$, so
		\[
		\lim_{h\to 0} \left[g(\Delta y)+ u'(y)\right] = g(0) + u'(y) = u'(y).
		\]
		Hence we have
		\[
		f'(x) = u'(v(x)) \cdot v'(x).
		\]
	\end{proofenv}
	\begin{remarkenv}
		In Leibniz's notation, suppose $y=v(x)$ and $z=u(y)$, then we write
		\[
		\frac{\dd y}{\dd x}:=v'(x)
		\quad \text{and}\quad
		\frac{\dd z}{\dd y}:=u'(y).
		\]
		We also write
		\[
		\frac{\dd z}{\dd x} := f'(x)
		\quad \text{where}\quad f=u \circ v.
		\]
		Then the chain rule can be written as
		\[
		\frac{\dd z}{\dd x} = \frac{\dd z}{\dd y} \cdot \frac{\dd y}{\dd x} \not\equiv \frac{\dd z}{\bcancel{\dd y}}\cdot \frac{\bcancel{\dd y}}{\dd x}.
		\]
		Here, $\not \equiv$ means ``not identically equal to", which indicates that the differentials $\dd y$ are \emph{not} actually quantities that can be cancelled out.
	\end{remarkenv}

	\section{Theorems of Derivatives}
	\begin{definitionenv}
		A function $f$ defined on a set $S \subseteq \mathbb{R}$ is said to have a \textbf{relative maximum} at $c\in S$ if there exists some open interval $I$ containing $c$ such that
		\[
		\forall x \in I \cap S,\ f(x) \le f(c).
		\]
		Similarly, we can define \textbf{relative minimum} of $f$ at $d\in S$.
	\end{definitionenv}
	\begin{remarkenv}
		We also say $f$ has an \textbf{extremum} at $c$ if it has either a relative maximum or a relative minimum at $c$.
	\end{remarkenv}
	\begin{theoremenv}
		Suppose function $f$ is defined on an open interval $(a,b)$ containing $c$. If $f$ has a relative extremum at $c$ and if $f'(c)$ exists, then
		\[
		f'(c)=0.
		\]
	\end{theoremenv}
	\begin{proofenv}
		Let
		\[
		g(x)=\begin{cases}
		\dfrac{f(x)-f(c)}{x-c}, & \text{if}\ x \in (a,b), x\neq c,\\[6pt]
		f'(c), & \text{if}\ x=c.
		\end{cases}
		\]
		It suffices to show that $g(c)=0$. Note
		\[
		\lim_{x \to c} \frac{f(x)-f(c)}{x-c}=f'(c)
		\]
		since $f'(c)$ exists, thus $g$ is continuous at $c$. Now suppose $g(c)>0$, then there exists some $\delta>0$ such that
		\[
		g(x)>0,\quad \forall x\in (c-\delta, c+\delta) \subseteq (a,b).
		\]
		If $x\in(c, c+\delta)$, then from $g(x)>0$ we derive $f(x)>f(c)$. If $x\in(c-\delta, c)$, then from $g(x)>0$ we derive $f(x)<f(c)$. Note this contradicts to $f$ having an extremum at $c$. Similarly, we can prove that $g(c)<0$ also leads to a contradiction. Thus we must have $g(c)=0$, i.e., $f'(c)=0$.
	\end{proofenv}
	\begin{theoremenv}[Rolle's Theorem]
		Let $f$ be continuous on $[a,b]$ and differentiable on $(a,b)$. If $f(a)=f(b)$, then there exists some $c\in (a,b)$ such that
		\[
		f'(c)=0.
		\]
		
	\end{theoremenv}
	\begin{proofenv}
		Suppose $f'(x)\neq 0$ for all $x\in (a,b)$. By \textbf{EVT}, we know
		\[
		\sup f = f(\alpha) \quad \text{and}\quad \inf f = f(\beta)\quad \text{for some}\quad \alpha, \beta \in [a,b].
		\]
		Note we know $f$ have extremum at $\alpha$ and $\beta$. If $\alpha, \beta \in (a,b)$, then by the above theorem we have
		\[
		f'(\alpha)=f'(\beta)=0,
		\]
		which contradicts our assumption. Thus $\alpha,\beta\in \{a,b\}$ and $f(\alpha)=f(\beta)$. Thus we know $f$ is constant on $[a,b]$, which implies
		\[
		f'(x)= 0,\quad \forall x\in (a,b),
		\]
		contradicting our assumption again. Then we complete the proof.
	\end{proofenv}
	\begin{theoremenv}[Lagrange's Mean Value Theorem]
		Let $f$ be continuous on $[a,b]$ and differentiable on $(a,b)$. Then there exists some $c\in (a,b)$ such that
		\[
		f'(c)=\frac{f(b)-f(a)}{b-a}.
		\]
		
	\end{theoremenv}
	\begin{proofenv}
		Define a new function
		\[
		g(x) = f(x) - \left[ \frac{f(b)-f(a)}{b-a} (x - a) + f(a) \right].
		\]
		Clearly $g$ is continuous on $[a,b]$ and differentiable on $(a,b)$. Note that
		\[
		\begin{cases}
			g(a)=f(a)-f(a)=0 ,\\[6pt]
			g(b)=f(b)-\left[ \dfrac{f(b)-f(a)}{b-a} (b - a) + f(a) \right]=0.
		\end{cases}
		\]
		By \textbf{Rolle's Theorem}, there exists some $c\in (a,b)$ such that
		\[
		g'(c) = f'(c) - \frac{f(b)-f(a)}{b-a} = 0,
		\]
		which implies
		\[
		f'(c) = \frac{f(b)-f(a)}{b-a}.
		\]
	\end{proofenv}
	\begin{theoremenv}[Cauchy's Mean Value Theorem]
		Suppose $f$ and $g$ are both continuous on $[a,b]$ and differentiable on $(a,b)$, then there exists some $c\in (a,b)$ such that
		\[
		f'(c)\left[ g(b)-g(a) \right] = g'(c) \left[ f(b)-f(a) \right].
		\]
		
	\end{theoremenv}
	\begin{proofenv}
		Define a new function
		\[
		h(x) = f(x) \left[ g(b)-g(a) \right] - g(x) \left[ f(b)-f(a) \right].
		\]
		Clearly $h$ is continuous on $[a,b]$ and differentiable on $(a,b)$. Note that
		\[
		\begin{cases}
			h(a)=f(a)\left[ g(b)-g(a) \right] - g(a)\left[ f(b)-f(a) \right],\\[6pt]
			h(b)=f(b)\left[ g(b)-g(a) \right] - g(b)\left[ f(b)-f(a) \right].
		\end{cases}
		\Rightarrow h(b)-h(a)=0.
		\]
		By \textbf{Rolle's Theorem}, there exists some $c\in (a,b)$ such that
		\[
		h'(c) = f'(c)\left[ g(b)-g(a) \right] - g'(c)\left[ f(b)-f(a) \right] = 0,
		\]
		which implies
		\[
		f'(c)\left[ g(b)-g(a) \right] = g'(c)\left[ f(b)-f(a) \right].
		\]
	\end{proofenv}
	\begin{theoremenv}
		Suppose $f$ is continuous on $[a,b]$ and differentiable on $(a,b)$. Then:
		\begin{enumerate}
			\item [(1)] If $f'(x)=0$ for all $x\in (a,b)$, then $f$ is constant on $[a,b]$;
			\item [(2)] If $f'(x)>0$ for all $x\in (a,b)$, then $f$ is strictly increasing on $[a,b]$;
			\item [(3)] If $f'(x)<0$ for all $x\in (a,b)$, then $f$ is strictly decreasing on $[a,b]$.
		\end{enumerate}
	\end{theoremenv}
	\begin{theoremenv}
		Suppose $f$ is continuous on $[a,b]$ and differentiable on $(a,b)$ and $c\in (a,b)$. Then:
		\begin{enumerate}
			\item [(1)] If $f'(x)>0$ for all $x\in (a,c)$ and $f'(x)<0$ for all $x\in (c,b)$, then $f$ has a relative maximum at $c$;
			\item [(2)] If $f'(x)<0$ for all $x\in (a,c)$ and $f'(x)>0$ for all $x\in (c,b)$, then $f$ has a relative minimum at $c$.
		\end{enumerate}
	\end{theoremenv}
	\begin{remarkenv}
		The above two cases still hold even if $f'(c)$ is not well-defined. (We still require $f$ to be continuous at $c$.)
	\end{remarkenv}
	\begin{definitionenv}
		Suppose $f$ is continuous on $[a,b]$ and differentiable on $(a,b)$ and $c\in (a,b)$. If $f'(c)=0$, then we call $c$ a \textbf{critical point} of $f$.
	\end{definitionenv}
	\begin{theoremenv}
		Suppose $f$ has a second-derivative $f''$ defined on $(a,b)$ and $c\in (a,b)$ is a critical point of $f$. Then:
		\begin{enumerate}
			\item [(1)] If $f''(c)>0$, then $f$ has a relative minimum at $c$;
			\item [(2)] If $f''(c)<0$, then $f$ has a relative maximum at $c$;
			\item [(3)] If $f''(c)=0$, then the test is inconclusive, i.e., $f$ may have a relative maximum, a relative minimum or neither at $c$.
		\end{enumerate}
	\end{theoremenv}
	\begin{theoremenv}
		Suppose $f$ is continuous on $[a,b]$ and differentiable on $(a,b)$. Then:
		\begin{enumerate}
			\item If $f'$ is increasing on $(a,b)$, or equivalently $f''(x) \geq 0$ for all $x \in (a,b)$ provided that $f''$ is well-defined, then $f$ is convex on $[a,b]$.
			\item If $f'$ is decreasing on $(a,b)$, or equivalently $f''(x) \leq 0$ for all $x \in (a,b)$ provided that $f''$ is well-defined, then $f$ is concave on $[a,b]$.
		\end{enumerate}
	\end{theoremenv}
	\begin{remarkenv}
		Note that a flat line is both convex and concave (in our text book's definition).
	\end{remarkenv}

	\newpage